\chapter[Model adequacy and physical resources]{Model adequacy, abstention, and physical resource guarantees}
\label{adequacy:chapter}

A decision rule that must choose between two mechanisms will still choose one
when both mechanisms are wrong. A scientific workflow therefore needs an
adequacy layer and an explicit abstention output. It also needs a resource ledger:
knowing how many operations are executed does not determine how much energy the
apparatus supplies. This chapter gives contracts under which these two layers
combine with the precision results of Chapter~\ref{precision:chapter}.

The probability bounds are mathematical guarantees about the stated predictive
and physical premises. They do not validate a detector, establish local detailed
balance, measure a fuel channel, or prove that the candidate family contains the
true system. Those premises must be independently justified for an application.

\section{Conditional predictions are the adequacy contract}

At round $t$, let $\mathcal F_{t-1}$ contain the full retained past, calibration
information used for this prediction, the next action, and any controller
randomization used to choose it. This is pre-observation information; one may
formally enlarge the filtration immediately before each outcome. A binary
statistic $Y_t$ is then revealed. The statistic and event being recorded must
also be chosen before seeing that outcome.

A candidate $j$ supplies predictable envelopes
$[\ell_{j,t},u_{j,t}]\subset[0,1]$. Under a fixed true model, conditional validity
means
\begin{equation}
 \ell_{j,t}\leq p_t:=\mathbb E[Y_t\mid\mathcal F_{t-1}]
       \leq u_{j,t}\quad\hbox{for every }t.
 \label{adequacy:conditional}
\end{equation}
For a composite candidate, the contract must hold under each fixed member for
which the candidate is asserted valid. Maximizing and minimizing correct
conditional forecasts over nuisance values can produce valid but conservative
envelopes. A fixed nuisance does not become independently redrawn at every round.

This contract allows dependent records, feedback and optional stopping. Its
strength is also its principal burden: a stationary marginal probability or
an unconditional error bound is not generally a conditional forecast. For a
simple failure control, let one fair coin $U$ be drawn and set $Y_t=U$ forever.
Every marginal is Bernoulli$(1/2)$. After the first observation, however, the
next conditional probability is zero or one. Treating these outcomes as fresh
fair coins can generate overwhelming, invalid rejection evidence against the
true marginal description.

\section{A sequential adequacy certificate}

Before $Y_t$ is revealed, choose a finite predictable stake $\lambda_{j,t}$.
For a nonnegative stake use $b_{j,t}=u_{j,t}$, and for a negative stake use
$b_{j,t}=\ell_{j,t}$. Define
\begin{equation}
 f_{j,t}(Y_t)=1+\lambda_{j,t}(Y_t-b_{j,t}),
 \qquad M_{j,n}=\prod_{t=1}^n f_{j,t}(Y_t).
 \label{adequacy:bet}
\end{equation}
Require the two factors at $Y_t=0$ and $Y_t=1$ to be nonnegative. For interior
endpoints this is $\lambda\leq1/u$ for positive stakes and
$\lambda\geq-1/(1-\ell)$ for negative stakes; checking both factors directly
handles zero and one endpoints without dividing by zero.

\begin{proposition}[Anytime exclusion]
\label{adequacy:anytime}
Under \eqref{adequacy:conditional}, $M_{j,n}$ is a nonnegative test
supermartingale starting at one. A fixed convex mixture $E_{j,n}$ of such
strategies has the same property. For any $0<\alpha_j<1$,
\begin{equation}
 \Pr\{\text{some }n:\ E_{j,n}\geq1/\alpha_j\}\leq\alpha_j.
 \label{adequacy:ville}
\end{equation}
\end{proposition}
\begin{proof}
The signed endpoint choice gives
\[
 \mathbb E[f_{j,t}(Y_t)\mid\mathcal F_{t-1}]
 =1+\lambda_{j,t}(p_t-b_{j,t})\leq1.
\]
Multiplying by past nonnegative wealth gives the supermartingale property;
fixed convex combinations preserve it. For completeness, stop at the first
threshold crossing or a deterministic horizon $n$. Optional stopping for this
bounded stopping time bounds expected stopped wealth by one. On a crossing,
wealth is at least $1/\alpha_j$, so its probability is at most $\alpha_j$.
Letting the horizon increase proves \eqref{adequacy:ville}.
\end{proof}

This is Ville's maximal inequality and bounded betting
\cite{caVille1939,caWaudbySmithRamdas2024}. Component stakes may depend on the
past. Fixed mixture weights are chosen before the tested record; replacing them
after seeing the current outcome requires a separate wealth-accounting argument.
The test is valid with any permitted predictable intervention policy satisfying
the envelopes. Validity alone says nothing about power or acquisition efficiency.

Exclude a candidate permanently once its threshold is crossed. Denote the
surviving set by $S_n$. Permanent exclusion records an ever-crossing event;
instantaneous wealth can fall again and does not undo that event.

\begin{theorem}[Target commitment with abstention]
\label{adequacy:commit}
Suppose one fixed candidate $j_*$ is valid throughout the run, contains the true
physical model, and has the true target label $g_*$. Each candidate is
target-homogeneous, with declared label $g_j$. Use level $\alpha_j=\alpha$
for every candidate. Return:
\begin{enumerate}
 \item \texttt{abstain\_all\_rejected} if $S_n$ is empty;
 \item \texttt{abstain\_ambiguous} if survivors assert different targets;
 \item \texttt{commit}$(v)$ if $S_n$ is nonempty and every survivor asserts $v$.
\end{enumerate}
The probability of ever returning an empty survivor set or ever committing an
incorrect target is at most $\alpha$.
\end{theorem}
\begin{proof}
If $j_*$ has not been excluded, the survivor set is nonempty and any unanimous
target equals $g_*$. Thus either unsafe event entails exclusion of $j_*$.
Proposition~\ref{adequacy:anytime} bounds that ever-exclusion probability.
\end{proof}

There is no factor equal to the number of candidates in this statement: it
protects one fixed valid and correctly targeted candidate and the resulting
decision. To protect every valid candidate simultaneously, use levels with
$\sum_j\alpha_j\leq\alpha$ and apply a union bound. These are different
questions. A different candidate fitting each successive prediction is not a
single candidate valid throughout the run.

If a candidate has unresolved internal target values, retain its target set
and return the union over survivors. Commitment to one value is justified only
when that union is a singleton. Observable agreement alone does not validate an
asserted hidden target: the theorem needs the true model to belong to a
correctly specified candidate. If truth is outside the declared family, both
wrong commitment and total rejection can occur without this guarantee.
Nonrejection is not evidence that all omitted mechanisms have been ruled out.
There is no general termination theorem; opposite-target models with identical
accessible laws may remain ambiguous indefinitely.

\section{Calibration uncertainty without invalid conditioning}

A joint calibration statement may guarantee that all conditional prediction
envelopes used during acquisition are valid with probability at least
$1-\delta$. One cannot usually prove an anytime guarantee by conditioning on
this eventual good event: that conditioning can change the conditional outcome
laws on which the betting calculation depends.

\begin{proposition}[Calibration composition]
\label{adequacy:calibration}
Let $G$ be any event on which all envelopes of the designated true candidate
are conditionally valid. If $\Pr(G^c)\leq\delta$, the unsafe-output probability
in Theorem~\ref{adequacy:commit} is at most $\alpha+\delta$, capped by one.
Moreover the probability of unsafe output intersected with $G$ is at most
$\alpha$.
\end{proposition}
\begin{proof}
The actual conditional mean $p_t$ and the supplied endpoints are predictable,
even if $p_t$ is unknown to the scientist. Stop the designated candidate's
wealth before the first round at which its envelope fails, replacing that and
all later factors by one. The stopped wealth is a nonnegative test
supermartingale. On $G$ it equals the original wealth throughout the run.
Unsafe output on $G$ therefore entails a crossing by this valid stopped wealth,
which has probability at most $\alpha$. Add $\Pr(G^c)$.
\end{proof}

A separate 95\% interval at every round does not provide one 5\% failure budget
for the experiment. Calibration must cover every conditional prediction used,
including drift, filtering, actions and detector behavior. An interface that
requires \texttt{prepare()} before \texttt{observe()} helps enforce chronology
but cannot prove that its caller did not inspect future data. Scientific record
retention remains part of the contract. Adding candidates invented after seeing
the tested data also requires new valid tests or a justified data-reuse scheme.

\section{What a physical energy budget must cover}

An additive action charge is a bookkeeping variable until calibrated to a
physical quantity. Time, supplied work, dissipated heat, net free-energy change,
and operation count are different quantities. Their units and scope must be
specified. In particular, a reversible added transition channel at a stationary
distribution does not make the apparatus or its controller free of energy cost.

For a complete experiment, the ledger must cover preparation, acquisition,
interventions, sensing and readout, controller operation, reset, intervening
gaps, and calibration. Shared support can be reported separately or amortized
over a declared number of experiments. Any omitted contribution prevents a
claim about total energy. A detector-law error bound does not bound energy in
unobserved physical channels.

Two meanings of a budget are particularly different. A pathwise bound is never
exceeded within its assumptions. An expected expenditure permits arbitrarily
large realized expenditure. A finite-time continuous-time Markov process can
make arbitrarily many jumps, so multiplying an expected event count by an
event energy does not produce a pathwise cap. The next theorem supplies a
probabilistic cap under a channel-level physical contract.

\section{A counting-channel acquisition cap}

Let $N_t$ count supplied-energy events, starting at zero. Assume it is a simple,
unit-jump, nonexplosive counting process with predictable intensity $\lambda_t$
and compensator
$A_t=\int_0^t\lambda_s\,ds$. Let $\tau$ be a bounded stopping time for a
filtration in which this intensity statement is valid. These are assumptions
on the actual process; observing $N_t$ or knowing its exact intensity online is
not required for the probability argument. A calibrated predictable intensity
upper bound may replace $\lambda_t$ and yields the corresponding
supermartingale bound.

\begin{theorem}[A cap valid at adaptive stopping]
\label{adequacy:count-tail}
For $\mu>0$ and an integer $n>\mu$,
\begin{equation}
 \Pr\{N_\tau\geq n,\ A_\tau\leq\mu\}
 \leq \exp\{-n\log(n/\mu)+n-\mu\}.
 \label{adequacy:count-tail-equation}
\end{equation}
In particular this bounds $\Pr(N_\tau\geq n)$ whenever
$A_\tau\leq\mu$ almost surely. If $A_\tau=0$ almost surely, then
$N_\tau=0$ almost surely.
\end{theorem}
\begin{proof}
For $s>0$, the counting-process stochastic exponential
\[
 Z_t=\exp\{sN_t-(e^s-1)A_t\}
\]
is a nonnegative local martingale and hence a supermartingale. This follows,
for instance, by applying the jump product rule to obtain
$dZ_t=Z_{t-}(e^s-1)(dN_t-\lambda_t\,dt)$ and then localizing
\cite{caBremaud1981}. At the bounded stopping time,
$\mathbb E Z_\tau\leq1$. On the event in
\eqref{adequacy:count-tail-equation},
$Z_\tau\geq\exp\{sn-(e^s-1)\mu\}$. Markov's inequality bounds its
probability by $\exp\{-sn+(e^s-1)\mu\}$. Minimize at
$s=\log(n/\mu)>0$. When $\mu=0$, use the same event bound for each $s>0$
and let $s$ increase to infinity for every positive integer $n$.
\end{proof}

The argument does not assume independent event increments or a Poisson count.
It does require the counting-process contract. A signed transfer matrix used
to compute output probabilities does not automatically supply such a physical
counting instrument.

Choose $n$ so that the right side of
\eqref{adequacy:count-tail-equation} is at most $\beta$. Suppose the gross work
per counted event is at most $w_{\max}\geq0$, and a separately justified bound
$W_{\rm app}$ covers all other supplied energy. On the event where this ledger
is valid,
\begin{equation}
 W_{\rm gross}\leq w_{\max}N_\tau+W_{\rm app}.
 \label{adequacy:energy-ledger}
\end{equation}
Consequently $W_{\rm cap}=nw_{\max}+W_{\rm app}$ is a conservative
$1-\beta$ energy cap when the compensator and ledger premises hold surely.
It is deliberately conservative by one count and is neither an exact count
quantile nor a hard pathwise cap. Signed net work must not replace gross
supplied work in \eqref{adequacy:energy-ledger} without an appropriate new model.

For rational $\mu$, $\beta$ and a proposed $n$, the cap can be checked without
trusting floating-point logarithms. For rational $x\geq1$, write
$z=(x-1)/(x+1)$ and use
\[
 L_K(x)=2\sum_{k=0}^{K-1}\frac{z^{2k+1}}{2k+1},\qquad
 L_K(x)\leq\log x\leq L_K(x)+
 \frac{2z^{2K+1}}{(2K+1)(1-z^2)}=:U_K(x).
\]
The positive series for $\log((1+z)/(1-z))$ proves the lower bound; replacing
each remaining denominator by $2K+1$ gives the geometric upper tail. Range
reduction by powers of two improves convergence. A rational certificate is
\[
 nL_K(n/\mu)-n+\mu\geq U_K(1/\beta).
\]
An independent checker recomputes these inequalities. Their truth validates the
probability calculation, while the laboratory remains responsible for the
physical intensity, event energy and apparatus bounds.

\section{A joint inference and energy guarantee}

\begin{theorem}[Composition without independence]
\label{adequacy:joint}
Assume one fixed correctly targeted candidate is present. Let $G$ be a joint
good event of probability at least $1-\delta$ on which:
\begin{enumerate}
 \item its full-history conditional prediction envelopes are valid throughout;
 \item $A_\tau\leq\mu$ for the actual counting process;
 \item the gross-energy ledger \eqref{adequacy:energy-ledger} holds.
\end{enumerate}
Run the target-commitment rule at level $\alpha$, and choose $n$ in
Theorem~\ref{adequacy:count-tail} at level $\beta$. Then
\begin{equation}
 \Pr\{\text{ever incorrect commitment, ever all-candidate rejection,
       or }W_{\rm gross}>W_{\rm cap}\}
 \leq\alpha+\beta+\delta,
 \label{adequacy:joint-equation}
\end{equation}
where inference outputs are considered through the bounded experimental stopping
time. No independence of calibration, inference and energy is required.
\end{theorem}
\begin{proof}
Unsafe inference intersected with $G$ has probability at most $\alpha$ by the
stopped-before-invalid-forecast argument in
Proposition~\ref{adequacy:calibration}. Energy overrun intersected with $G$
implies $N_\tau\geq n$ and $A_\tau\leq\mu$, unless $w_{\max}=0$, in which case
it is impossible on $G$. Theorem~\ref{adequacy:count-tail} bounds that event by
$\beta$. Add the two events and $G^c$ by the union bound.
\end{proof}

The compensator or energy-ledger failure need not be predictable: their terminal
joint events are used directly. In contrast, conditional forecasts are stopped
before their first invalid factor to preserve the supermartingale. This
distinction avoids an unjustified application of optional stopping or
conditioning. A single joint calibration failure budget is counted once.

\section{An end-to-end handoff from a witness to a physical guarantee}

The following result gives a common handoff between the algebraic, statistical
and physical layers. It uses a certified score gap; the exact witness that
motivated that score need not itself optimize the final energy cost.

\begin{corollary}[Robust score, record count, and energy cap]
\label{adequacy:handoff}
Fix target classes $+$ and $-$, a predeclared protocol and bounded record score,
and choose constants $0\leq b<a\leq1$ with $\gamma=a-b>0$.
Let $M$ be a fixed integer with
$M\gamma^2/2\geq\log(1/\alpha)$. Assume an event $G$ of probability at least
$1-\delta$ simultaneously ensures:
\begin{enumerate}
 \item under the true fixed target class, every actual score's conditional
 mean is at least $a$ for $+$ or at most $b$ for $-$;
 \item the complete bounded acquisition has $A_\tau\leq\mu$ and satisfies
 the gross-energy ledger \eqref{adequacy:energy-ledger}.
\end{enumerate}
The midpoint score decision after $M$ records, with any count cap $n$ certified
at level $\beta$, obeys
\begin{equation}
 \Pr\{\text{wrong target or }W_{\rm gross}>nw_{\max}+W_{\rm app}\}
 \leq\alpha+\beta+\delta
 \label{adequacy:handoff-equation}
\end{equation}
under every fixed model satisfying those contracts.
\end{corollary}
\begin{proof}
On $G$, the conditional score means obey the required one-sided bound. To avoid
conditioning on $G$, stop at the first predictable violation of that bound and
replace that and all subsequent scores by $a$ under $+$, or by $b$ under $-$.
The modified sequence has the required conditional means everywhere and agrees
with the executed sequence on $G$. The conditional version of
Proposition~\ref{precision:score-test} therefore bounds wrong decision
intersected with $G$ by $\alpha$. Bound energy overrun intersected with $G$
by Theorem~\ref{adequacy:count-tail}, and add $\Pr(G^c)$.
\end{proof}

Independent identically prepared records from a fixed actual model are a
sufficient way to obtain the score contract. The conditional formulation also
states precisely what would justify using dependent records. If nominal
expectation bounds and total-variation calibration produce the gap, use
Equation~\eqref{precision:robust-gap}; do not amplify an uncalibrated nominal
contrast and add one-record error only at the end. If fresh reset fails or an
intervention changes the nuisance between records, the model must be rebuilt.

For a simple synthetic calculation, let $a=3/5$, $b=2/5$, and $\alpha=1/20$.
Then $\gamma=1/5$ and $M=150$ certifies each conditional decision error because
$M\gamma^2/2=3>\log20$. Suppose the separately specified complete physical
contract has compensator cap $\mu=150$. The rational counting checker certifies
$n=189$ at $\beta=1/100$. With joint calibration failure $\delta=1/100$,
the probability of a wrong decision or exceeding
$189w_{\max}+W_{\rm app}$ is at most $7/100$.
These numbers are an arithmetic demonstration, not measured performance.
If native fuel use continues during preparation or gaps, that exposure must be
included in $\mu$ or bounded separately in the ledger. The declaration
$\mu=150$ is insufficient until all relevant physical time has been accounted for.

\section{What a scientist should retain}

A usable result consists of the target and candidate definition; the actual
observation, preparation and intervention contract; an exact obstruction or
witness where available; quantitative calibration and sampling bounds; the
executed policy and complete records; and the physical resource ledger with
its coverage statement. The inference should return ambiguity or failed model
adequacy when those are the supported conclusions. A finite certificate is
portable because another scientist can check its mathematical content. Its
application becomes credible when another apparatus independently verifies
the physical premises and reproduces the predicted distinction.
